\textbf{LLM/VLM planning for robotics.} Foundation models have been used to
ground natural-language instructions in robotic
affordances~\cite{ahn2022saycan}, to act as zero-shot
planners~\cite{huang2022zeroshot}, and to generate feasible motion plans
directly from instructions~\cite{lin2023text2motion}. A range of reasoning
strategies has been proposed on top of these capabilities: reactive,
single-step prompting; interleaved reasoning-and-acting
(ReAct)~\cite{yao2022react}; iterative self-critique
(Self-Refine)~\cite{madaan2023selfrefine}, including a long-horizon
sequential-planning variant validated at ICRA~\cite{zhou2024isrllm};
search-based deliberation (Tree-of-Thoughts)~\cite{yao2023tot}; and
ensemble voting over independently sampled plans
(chain-of-thought with self-consistency)~\cite{wei2022cot,
wang2022selfconsistency}. Recent surveys catalogue this space in
detail~\cite{kim2024survey, wang2025robotics, huang2024planningsurvey,
berti2025emergent, billard2025roadmap}. Our contribution is orthogonal to
this line of work: we do not propose a new reasoning strategy, and hold
the reasoning strategy fixed in our evaluation (\S\ref{sec:evaluation}).
Instead, we show that the geometric grounding surrounding \emph{any} of
these strategies -- how a chosen action's numeric parameters are computed
-- is a separate failure axis that persists regardless of which reasoning
strategy produced the action.

\textbf{End-to-end and vision-language-action approaches.} An alternative
architecture collapses perception, reasoning, and control into a single
model, as in vision-language-action (VLA) systems~\cite{brohan2023rt2,
driess2023palme} or code-generating policies that couple perception
directly to control primitives~\cite{liang2023codeaspolicies,
huang2023voxposer, huang2022innermonologue, singh2023progprompt}. These
approaches can, in principle, learn to compensate for geometric grounding
errors end-to-end, at the cost of the diagnosability that a decoupled
Sense-Plan-Act pipeline provides -- when an end-to-end model fails, it is
difficult to determine whether the failure was a misperception, a
reasoning error, or a motor-control error~\cite{favaliRAS}. We deliberately
retain a decoupled architecture, precisely because it lets us isolate and
name the geometry-grounding failure class this paper documents; a fully
end-to-end system would not permit the diagnosis in \S\ref{sec:taxonomy}
at all, only its symptoms.

\textbf{Simulation and sim-to-real transfer.} The sim-to-real gap is a
long-studied concern in robot learning, primarily for learned low-level
control policies~\cite{zhao2020simtoreal}. Physics simulators such as
Gazebo~\cite{koenig2004gazebo}, MuJoCo~\cite{todorov2012mujoco}, and
CoppeliaSim~\cite{rohmer2013coppeliasim} narrow this gap for control by
modeling contact dynamics and actuation limits~\cite{collins2021physics}.
The gap we document is different in kind: it does not concern whether a
\emph{learned policy} generalizes from simulated to real dynamics, but
whether a \emph{language-model-computed} geometric quantity -- one that a
symbolic benchmark never forces the model to compute against a
contradicting sensor -- is trustworthy at all. To our knowledge this
failure class has not been previously isolated and named for LLM/VLM robot
planners specifically.

\textbf{Direct comparison.} The closest prior work to this paper is
Favali et al.~\cite{favaliRAS}, which introduces a task/environment
complexity taxonomy (task length, goal specificity, affordance
perception; affordance density and four classes of dynamic disturbance)
and TS/TSR/AETS metrics for benchmarking LLM-based robotic planning
across six reasoning strategies, evaluated in a symbolic simulator with
exact ground truth. We adopt the same evaluation philosophy and, for
\S\ref{sec:evaluation}, the same metric family, but on real hardware,
where ground truth for safety and success is not automatically available
(\S\ref{sec:evaluation} details the resulting semi-automatic
adaptation). A companion systems paper from the same group,
RoboReason~Lab~\cite{favaliRoboReasonLab}, provides a modular simulation
testbed for the same class of reasoning strategies in CoppeliaSim; like
\cite{favaliRAS}, it never grounds a plan in measured depth or a
calibrated camera. Both are simulation-only by design, and both are
structurally unable to produce the geometry-grounding failures this paper
diagnoses, because neither ever asks a model to compute a quantity a
physical sensor could contradict.
